problem:  
Let \((M,\mathcal{D},g)\) be a **3‑dimensional contact sub‑Riemannian manifold** equipped with the canonical Popp measure and sub‑Laplacian \(\Delta\). Let \(p_t(x,y)\) denote the corresponding heat kernel. The manifold has two primary scalar curvature invariants: the **Tanaka–Webster scalar curvature** \(\kappa(x)\) and the **torsion invariant** \(\chi(x)\), both entering asymptotic expansions and sub‑Riemannian curvature formulas (Agrachev–Barilari–Paoli).  

It is known that, as \(t\to 0^+\),
\[
p_t(x,x)\sim \frac{C}{t^{Q/2}}\Bigl(1+\beta_1(x)t+\beta_2(x)t^2+\cdots\Bigr),\quad Q=4,
\]
where \(\beta_j(x)\) depend on the local curvature invariants and their derivatives. For the Heisenberg group \(H^1\), all curvature invariants vanish and \(\beta_j(H^1)=0\) for all \(j\ge1\).  

**Open problem:**  
In the class of compact, equiregular 3‑dimensional contact sub‑Riemannian manifolds, assume that the first two nontrivial heat kernel coefficients satisfy
\[
\beta_1(x)\equiv 0,\qquad \beta_2(x)\equiv 0 \ \ \text{for all } x\in M.
\]
Does it follow that both curvature invariants \(\kappa(x)\) and \(\chi(x)\) vanish identically, and hence that \((M,\mathcal{D},g)\) is locally isometric to the Heisenberg group \(H^1\)?  
Equivalently, do the first few on‑diagonal coefficients of the heat kernel determine local flatness of a 3D contact sub‑Riemannian structure?

Why is it a "good" problem:  
This formulation avoids ambiguity by restricting to the simplest nontrivial equiregular case—3D contact geometry—where explicit curvature invariants and partial asymptotic expansions are known. The equality of early heat coefficients with those for the flat model (the Heisenberg group) offers a concrete analytic condition whose geometric implications are unknown. Resolving it would precisely quantify how much geometric information is encoded in the beginning of the heat expansion, a foundational question bridging **microlocal analysis** (small‑time heat kernel asymptotics) and **differential geometry** (contact and sub‑Riemannian curvature).  

The statement is simple—“Do the first two heat coefficients determine flatness?”—yet profoundly deep because it asks whether the sub‑Riemannian curvature tensor can be recovered from finite spectral data. This connects spectral geometry, PDE theory of hypoelliptic operators, and the intrinsic geometry of non‑holonomic distributions. Its answer, positive or negative, would sharpen our understanding of geometric rigidity and spectral invariants in the sub‑Riemannian realm, exemplifying the qualities of a “good” research problem.